% Wampanoag (Trumbull 1903) — Source Workflow
% Issue #1363
% Status: Research complete, implementation pending

\section{Wampanoag: Trumbull/Eliot 1903}
\label{sec:1363}

\subsection{Source Description}

\textbf{Recorder:} James Hammond Trumbull (1821--1897), drawing on John Eliot's
17th-century Massachusett writing system.
\textbf{Source:} \emph{Natick Dictionary}, Bureau of American Ethnology
Bulletin 25 (1903, posthumous).
\textbf{Language:} Wampanoag (Massachusett) [wam].
\textbf{Corpus:} 4,491 records (largest in the SNEA collection).
\textbf{Orthography:} Eliot orthography --- the first standardized writing
system for an Algonquian language, developed by John Eliot and Massachusett
collaborators (Job Nesutan, John Sassamon) in the 1650s--1660s.

\subsubsection{Recorder Biography: James Hammond Trumbull}

James Hammond Trumbull (1821--1897) was an American philologist, historian,
and bibliographer. He served as Connecticut State Librarian (1854--1897) and
was a founding member of the American Philological Association. His work on
the Algonquian languages of southern New England spanned four decades,
beginning with his early comparative study \emph{On the Best Method of
Studying the North American Languages} (1859).

Trumbull's major Algonquianist works include:
\begin{itemize}
  \item \emph{The Composition of Indian Geographical Names} (1870) ---
    Connecticut Historical Society
  \item \emph{Indian Names of Places in Connecticut} (1881) --- systematic
    geographic toponymy
  \item \emph{Natick Dictionary} (published posthumously 1903) --- his
    magnum opus
\end{itemize}

\subsubsection{The Lost Thesis Controversy}

Trumbull delivered a landmark paper before the American Philological
Association titled \emph{The Method of the Indian Languages as Exemplified by
the Massachusett Dialect} --- based on a ``lost thesis'' he had written at
Yale in the 1840s. The paper argued for the polysynthetic structure of
Algonquian languages and established Trumbull's reputation as the leading
American philologist of Native languages before the Boasian revolution.

Trumbull's methodology combined:
\begin{enumerate}
  \item \textbf{Text-based extraction:} Exhaustively reading through all of
    Eliot's Massachusett translations (Bible, primers, catechisms) and
    manually extracting every lexical item onto slips
  \item \textbf{Cross-referencing with cognates:} Comparing Massachusett
    forms with related languages (Narragansett via Williams 1643,
    Mohegan-Pequot via field notes, Abenaki, Micmac, etc.)
  \item \textbf{Morphemic analysis:} Breaking down polysynthetic forms into
    constituent morphemes --- a radical approach for its time
  \item \textbf{Etymological comparison:} Identifying cognate sets across the
    Eastern Algonquian languages
\end{enumerate}

\subsubsection{Posthumous Publication}

Trumbull died in 1897 with the dictionary manuscript substantially complete
but not yet in press. Edward Everett Hale (the Unitarian minister and author)
supervised the final editing at the Bureau of American Ethnology. The preface
acknowledges that some inconsistencies in cross-references and source
abbreviations may reflect the incomplete final revision by Trumbull himself.
The published dictionary includes 4,491 main entries.

\subsubsection{The Eliot Orthography --- Historical Origins}

John Eliot (1604--1690) arrived in Massachusetts Bay in 1631 and began
preaching to the Massachusett people in the 1640s. His language teacher was
\textbf{Job Nesutan}, a Massachusett man who served as Eliot's primary
linguistic informant and translator. Eliot also worked with \textbf{John
Sassamon} (a Massachusett convert educated at Harvard) who later became a
crucial intermediary.

Eliot's orthography was developed in the 1650s as he prepared his translation
of the Bible. It was:
\begin{itemize}
  \item \textbf{Phonemically oriented} for its time --- Eliot attempted to
    represent the sounds of Massachusett systematically, not just through
    English ear-approximations
  \item \textbf{English-based} --- using English letter values as the starting
    point, then extending them
  \item \textbf{Regularized} --- unlike Wood (1634) or Winslow (1624), Eliot
    chose consistent spellings for the same morpheme across occurrences
  \item \textbf{Influenced by Hebrew and Greek grammar} --- Eliot's
    \emph{Indian Grammar Begun} (1666) explicitly draws parallels between
    Massachusett structures and Biblical Hebrew
\end{itemize}

The \emph{Mamusse Wunneetupanatamwe Up-Biblum God} (1,180 pages) was printed
by Samuel Green and Marmaduke Johnson in Cambridge, Massachusetts. It was the
first complete Bible printed in the Americas in a Native language. Production
required a specially cast type font (some characters unique to Massachusett),
translators and typesetters literate in the orthography, and Massachusett
readers to proofread.

\subsection{Vowel Inventory}

\subsubsection{Simple Vowels}

The core challenge of the Eliot orthography is that 5 vowel letters must
represent a larger Massachusett vowel inventory (likely 8--10 vowel phonemes
including length distinctions).

\begin{table}[h]
\centering
\begin{tabular}{lllll}
\toprule
Eliot & Phoneme & IPA & Example & Notes \\
\midrule
a & Short a / Schwa & /a/ ~ /ə/ & \textit{wadchu} (mountain) & Unstressed → [ə] \\
a & Long a & /aː/ & \textit{mahche} (already) & Lengthened \\
e & Long e & /iː/ & \textit{wetu} (house) & Two different underlying vowels \\
e & Short e & /ɛ/ & \textit{penashimwog} & \\
e & Final schwa & /ə/ & \textit{nish} (two) & Often silent but written \\
i & Long i & /iː/ & \textit{pish} (future) & \\
i & Short i & /ɪ/ & \textit{nippe} (water) & \\
o & Long o & /uː/ & \textit{noh} (father) & \\
o & Short o & /ɔ/ & \textit{nuppa} (sleep) & \\
o & Schwa & /ə/ & \textit{onk} (and) & Unstressed reduction \\
u & Schwa & /ə/ & \textit{ummon} (breast) & Unstressed position \\
u & Short u & /ʌ/ & \textit{musqui} (grass) & \\
\bottomrule
\end{tabular}
\caption{Simple vowel mappings in Eliot's orthography}
\label{tab:1363-vowels}
\end{table}

\subsubsection{Vowel Digraphs}

\begin{table}[h]
\centering
\begin{tabular}{llll}
\toprule
Eliot Digraph & IPA & Example & Notes \\
\midrule
ee & /iː/ & \textit{neen} (I), \textit{keetompauog} (men) & Consistent --- most reliable \\
oo & /uː/ & \textit{woskot} (firewood), \textit{sook} (pouring) & Consistent \\
⟨ꝏ⟩ / ∞ & /uː/ or complex & See §\ref{sec:1363-infinity} & The debated symbol \\
aa & /aː/ & \textit{waam} (long), \textit{maangu} (wolf) & \\
ii & /iː/ & \textit{niin} (my) & Less common \\
au & /ɔ/ & \textit{mauche} (come) & Rare \\
ai & /aːj/ & \textit{wai} (woe) & Diphthong \\
eu & /eːw/ & \textit{seup} (river) & Possibly /əw/ \\
\bottomrule
\end{tabular}
\caption{Vowel digraphs}
\label{tab:1363-digraphs}
\end{table}

\subsection{Consonant Inventory}

\begin{table}[h]
\centering
\begin{tabular}{llll}
\toprule
Eliot & IPA & Example & Notes \\
\midrule
p & /p/ ~ /b/ & \textit{peche} (tomorrow), \textit{nuppa} (sleep) & Unaspirated, lenis medial \\
t & /t/ ~ /d/ & \textit{tummunk} (worm), \textit{wetu} (house) & Unaspirated, medially lenis \\
k & /k/ ~ /ɡ/ & \textit{kenug} (fish), \textit{askkuhnk} (serpent) & Unaspirated \\
c (a/o/u) & /k/ & \textit{commaco} (feast), \textit{cutsh} (speak) & Hard c --- English convention \\
c (e/i) & /s/ & \textit{ce} (thou), \textit{nuppacoh} (chaff) & Soft c --- English convention \\
ch & /tʃ/ & \textit{choh} (child), \textit{mache} (go) & Affricate \\
sh & /ʃ/ & \textit{asuhsh} (fine), \textit{shuckquan} (rain) & \\
s & /s/ & \textit{simma} (various) & Word-initial s often written ss \\
ss & /sː/ ~ /s/ & \textit{massa} (great) & Geminate \\
m & /m/ & \textit{matta} (no) & \\
n & /n/ & \textit{neen} (I) & \\
w & /w/ & \textit{wetu} (house), \textit{waantam} (wise) & \\
y & /j/ & \textit{yeuyeu} (gradually) & \\
q & /kʷ/ & \textit{qussuk} (stone) & Labialized k --- from PA *kw \\
h & /h/ & \textit{ohke} (earth), \textit{kehtan} (great) & Also marks preaspiration \\
ht & /ht/ ~ /h/ & \textit{kehtan} (great), \textit{kehte} (truly) & Preaspirated /t/ \\
hp & /hp/ ~ /h/ & \textit{wuhpom} (honey) & Preaspirated /p/ \\
hk & /hk/ ~ /h/ & \textit{ohke} (earth) & Preaspirated /k/ \\
' (apostrophe) & /h/ ~ /ʔ/ & \textit{m'tah} (be tired) & Elided syllable, often = /h/ \\
\bottomrule
\end{tabular}
\caption{Consonant inventory}
\label{tab:1363-consonants}
\end{table}

\subsection{The ∞ Symbol --- Historical Typography and Phonetics}
\label{sec:1363-infinity}

\subsubsection{The Misidentification}

The ∞ symbol (Unicode INFINITY, U+221E) is one of the most distinctive
features of the database's Trumbull entries. However, this is almost
certainly a \textbf{typographic misidentification} in the digitization
process. The original Trumbull dictionary uses a different symbol:

The actual symbol: Eliot's Bible and Trumbull's dictionary use a ligature of
two ⟨o⟩ letters joined --- what typographers call an ``oo ligature'' ---
Unicode ⟨ꝏ⟩ (U+A74F, LATIN SMALL LETTER OO). In the early BAE Bulletin 25
printed type, this appeared as a figure-8-like symbol because the two o's
were cast as a single piece of type:

\begin{quote}
Original Eliot/Trumbull: ⟨ꝏ⟩ (oo ligature)\\
What it looks like at small sizes: ∞ (infinity-like)\\
What DB has: ∞ (Unicode INFINITY)
\end{quote}

The digitizer apparently did not recognize the typographic convention and
encoded it as the visually-similar infinity symbol.

\subsubsection{What ⟨ꝏ⟩ Actually Represents}

In Eliot's orthography, ⟨ꝏ⟩ represents \textbf{long /uː/} --- specifically,
it is Eliot's way of writing the long u sound, contrasting with ⟨oo⟩ which
may represent a different vowel quality or length. Eliot's own description in
\emph{The Indian Grammar Begun} (1666):

\begin{quote}
``oo, and ꝏ, are two differing vowels, or rather the same vowel differently
sounded; oo is short and ꝏ is long.''
\end{quote}

This suggests ⟨oo⟩ = /u/ (short) and ⟨ꝏ⟩ = /uː/ (long). However, database
evidence showing 654 occurrences of the symbol (now encoded as ∞) suggests a
more complex distribution. In many entries, ∞ appears where a
preaspirated/breathy vowel sequence would be expected phonologically.

\subsubsection{The ∞→oozzz Pipeline Mapping}

Per the lessons-learned documentation, the mapping ∞ → oozzz was established
empirically based on:
\begin{enumerate}
  \item \textbf{Distributional patterns:} ∞ appears in contexts where /u/ +
    /h/ + /z/ would be expected from PA reconstruction
  \item \textbf{Morphophonemic alternation:} Where cognate forms in other SNEA
    languages show /u/ + preaspiration + sibilant sequences
  \item \textbf{Goddard \& Bragdon (1988) transcription conventions:} Their
    modern phonemic transcription suggests underlying forms with /uːh/
    sequences
\end{enumerate}

The \texttt{oozzz} string is therefore a \textbf{working normalization} --- it
captures the phonetic complexity (length + breathiness + sibilant quality)
that the original ⟨ꝏ⟩ was designed to encode, even if the exact phonetics
remain debated.

\begin{table}[h]
\centering
\begin{tabular}{lll}
\toprule
Pattern & Count (approx.) & Example \\
\midrule
∞ (bare) & ~500 & \textit{achm∞wonk} (news) \\
∞̌ (with breve) & ~50 & \textit{∞̌ppa} (various) --- phonetically reduced \\
∞ in compounds & ~100 & Various morpheme-internal occurrences \\
\midrule
Total & ~654 & 14.6\% of 4,491 records \\
\bottomrule
\end{tabular}
\caption{∞ in the database corpus}
\label{tab:1363-infinity-counts}
\end{table}

\subsection{Entry Structure --- Detailed Parsing Guide}

Trumbull entries have a complex multi-part structure parsed into a single
\texttt{lx} field. Understanding this structure is essential for FSM-based
parsing.

\subsubsection{Standard Entry Template}

\begin{quote}
\texttt{*|achm∞wonk|, vbl. n. `news'    C.}
\end{quote}

Field breakdown:
\begin{itemize}
  \item \texttt{*} --- Cross-reference marker
  \item \texttt{|...|} --- Citation form delimiters (morphemic boundaries)
  \item \texttt{,} --- Field separator
  \item \texttt{vbl. n.} --- Grammatical label (verbal noun)
  \item \texttt{`news'} --- Gloss in backtick quotes
  \item \texttt{C.} --- Source abbreviation
\end{itemize}

\subsubsection{Entry Variants}

Root/compound form (no gloss):
\begin{quote}
\texttt{|-adchaubuk|, in comp. words, root, or roots}
\end{quote}

Cross-reference entry:
\begin{quote}
\texttt{|asqhutt∞che| `whilst'    C. = |asq-utt∞che|.}
\end{quote}

Multiple glosses:
\begin{quote}
\texttt{|wame|, adj. `all, whole, every'    C.}
\end{quote}

\subsubsection{Source Abbreviations}

Trumbull used a system of abbreviations to tag each entry's source:

\begin{table}[h]
\centering
\begin{tabular}{lll}
\toprule
Abbreviation & Source & Approx. \% \\
\midrule
C. & Eliot's Catechism (various editions) & ~40\% \\
B. & Eliot's Bible (1663, 1685) & ~30\% \\
P. & Psalms translated by Eliot & ~10\% \\
G. & Eliot's Grammar (1666) & ~5\% \\
Pr. & Eliot's Primer(s) & ~5\% \\
Pl. & Plymouth/colonial records & ~3\% \\
(unmarked) & Unknown or Trumbull's own reconstruction & ~7\% \\
\bottomrule
\end{tabular}
\caption{Source abbreviations in Trumbull}
\label{tab:1363-sources}
\end{table}

\subsubsection{Grammatical Labels}

\begin{table}[h]
\centering
\begin{tabular}{lll}
\toprule
Abbreviation & Full Form & Meaning \\
\midrule
vbl. n. & verbal noun & Action nominalization \\
adj. & adjective & Modifier \\
pro. & pronoun & Personal/demonstrative pronoun \\
adv. & adverb & Modifier of verb/adjective \\
prep. & preposition & Relational word \\
conj. & conjunction & Connective \\
interj. & interjection & Exclamation \\
n. & noun & Substantive \\
v. & verb & Verb (unmarked for transitivity) \\
v. a. & verb active & Transitive verb \\
v. n. & verb neuter & Intransitive verb \\
part. & particle & Uninflected word \\
s. & substantive & Noun (alternative to n.) \\
pl. & plural & Plural form \\
comp. & compound & Compound word \\
in comp. & in compound & Used only in compounds \\
dim. & diminutive & Diminutive form \\
\bottomrule
\end{tabular}
\caption{Grammatical label abbreviations}
\label{tab:1363-grammar}
\end{table}

\subsection{Problems Encountered}

\subsubsection{The ∞ (ꝏ) Symbol Misidentification}

The symbol was historically misidentified as an infinity symbol (∞). It is
Eliot's ⟨oo⟩ ligature (U+A74F), representing long /uː/. The digitization
error encoding it as Unicode INFINITY (U+221E) has propagated through the
database, requiring normalization in the conversion pipeline.

\subsubsection{Preaspiration Gap}

Eliot writes \textit{hk/hp/ht} only ~15\% of expected positions. Edwards
(1787 Mahican) is the only colonial source that consistently marks
preaspiration. For Trumbull/Eliot entries, preaspiration /h/ must be inferred
from comparative evidence (Edwards' Mahican data, PA reconstructions,
morphotactic position) rather than read directly from the orthography.

\subsubsection{Apostrophe Ambiguity}

The apostrophe may represent /h/, /ʔ/, or vowel elision --- no consistent
rule. The apostrophe function is ambiguous without comparative evidence and
cannot be resolved solely from within the Trumbull entry.

\subsubsection{Dialect Stratification}

Variation between Natick (Eliot) and Pokanoket (Winslow/Wood) may be
orthographic, phonological, or dialectal. The \textit{achm∞wonk} vs.
\textit{hobbamock} divergence (the ``devil problem'') suggests real dialect
differentiation between the ``praying town'' dialect of Natick and the
Plymouth-area Pokanoket dialect.

\subsubsection{Vowel Quality}

The five-letter vowel system (a, e, i, o, u) must represent 8--10 phonemes.
Disambiguation requires distributional analysis, PA correspondence, stress
position, and morphological position.

\subsubsection{\textit{c} Ambiguity}

Before back vowels (a, o, u) → /k/; before front vowels (e, i) → /s/.
Exception: some irregular forms where \textit{c} = /k/ before e. Word-final
\textit{c} → /k/.

\subsubsection{Boundary Markers}

The \texttt{|...|} pipe convention is near-universal (4,488 of 4,491 entries)
but embedded glosses inside pipes are rare, cross-reference syntax (=) may
span pipe boundaries, and hyphens within pipes indicate morpheme boundaries,
not word boundaries.

\subsubsection{The Devil Problem: \textit{achm∞wonk} vs. \textit{hobbamock}}

One of the most revealing lexical puzzles in the corpus involves the word for
`devil':

\begin{table}[h]
\centering
\begin{tabular}{lll}
\toprule
Source & Word & Notes \\
\midrule
Winslow (1624) & \textit{hobbamock} & Initial h written; also a proper name \\
Wood (1634) & \textit{abamocho}, \textit{abomocho} & No initial h; o-final \\
Williams (1643) & \textit{ábammocho} & Narragansett; consistent with Wood \\
Trumbull/Eliot & \textit{achm∞wonk} & Different root; \textit{ch} not \textit{b/p} \\
\bottomrule
\end{tabular}
\caption{The devil problem across sources}
\label{tab:1363-devil}
\end{table}

This shows that Trumbull/Eliot may be drawing from a different lexical
stratum (the ``praying town'' dialect of Natick) than Winslow/Wood
(Plymouth-area Pokanoket). The lexical difference is not just orthographic
--- it may reflect real dialect divergence.

\subsection{Research Approach}

\subsubsection{Recommended Methodology}

\textbf{DO NOT USE REGEX} --- the ∞ symbol, pipe boundaries, embedded
glosses, and cross-reference syntax require context-sensitive parsing.

\textbf{Stage 1: Preprocessing FSM.} Tokenize Trumbull entries by:
\begin{enumerate}
  \item Recognizing \texttt{|...|} citation boundaries and stripping pipes
  \item Identifying ∞ and replacing with \texttt{oozzz} (first normalization)
  \item Extracting backtick-delimited glosses
  \item Isolating grammatical abbreviations (vbl. n., C., etc.)
  \item Handling \texttt{*} prefix and \texttt{=} cross-reference syntax
  \item Separating source abbreviations (C., B., etc.) for provenance tracking
\end{enumerate}

\textbf{Stage 2: Phonemic Conversion FSM.} Apply Eliot-to-phoneme mapping
rules:
\begin{enumerate}
  \item Map Eliot vowel digraphs (ee, oo, aa, etc.) to phonemes
  \item Handle \textit{c}→\textit{k}/\textit{s} disambiguation based on
    following vowel
  \item Infer preaspiration /h/ from written \textit{hk}, \textit{hp},
    \textit{ht} clusters (direct evidence), morphotactic position (coda of
    stressed syllable, before /p t k/), comparative evidence from Edwards
    (1787) Mahican data, and PA reconstructions showing */h/ in corresponding
    position
  \item Resolve ∞→\texttt{oozzz} (already mapped in Stage 1)
  \item Resolve apostrophe to /h/ or /ʔ/ based on context and comparative
    evidence
\end{enumerate}

\textbf{Stage 3: Comparative Fallback.} For ambiguous vowel quality (the
5-vowel→8--10-phoneme mapping problem):
\begin{enumerate}
  \item PA comparison: Use Aubin (1975) reconstructions as primary
    disambiguation
  \item Sister-language evidence: Williams (Narragansett), Edwards (Mahican),
    Fielding/Prince-Speck (Mohegan-Pequot)
  \item Known cognate patterns: Eliot orthography ↔ modern Wampanoag
    revitalization (Weeden, Baird)
  \item Goddard \& Bragdon (1988) transcriptions: Modern phonemic analyses of
    Massachusett texts
\end{enumerate}

\textbf{Stage 4: Statistical/Distributional Analysis.} The large corpus (4,491
entries) enables statistical approaches not viable for smaller sources:
\begin{enumerate}
  \item Vowel quality prediction model based on surrounding consonant
    environment and stress position
  \item Preaspiration probability model based on morphotactic position and PA
    correspondence
  \item ∞ behavior model: Distributional patterns of ∞ across phonological
    environments
\end{enumerate}

\subsubsection{Proto-Algonquian Correspondences}

\begin{table}[h]
\centering
\begin{tabular}{llll}
\toprule
PA & Massachusett Reflex & Eliot Grapheme & Example \\
\midrule
*a & /a/ & a & *aʔlapy → \textit{anupp} (bowstring) \\
*e & /ə/ & u, a & *eːθkweːwa → \textit{squa} (woman) \\
*i & /i/ & i, e & *iːli- → \textit{en-} (thus) \\
*o & /ɔ/ & o, a & *oːθaːna → \textit{-adchan} \\
*aː & /aː/ & aa, a & *aːmoːwa → \textit{amaug} (bee) \\
*eː & /iː/ & ee & *eːhšwa → \textit{ash} (three) \\
*iː & /iː/ & ee & *iːniwa → \textit{neen} (I) \\
*oː & /uː/ & oo, ꝏ & *oːtaːna → \textit{adchan} \\
*p & /p/ & p & *pexpenya → \textit{pepeng} \\
*t & /t/ & t & *taːnya → \textit{ta} (there) \\
*k & /k/ & k, c & *kaːkinoː → \textit{kakine} (all) \\
*m & /m/ & m & *maškyekwi → \textit{maskiht} (medicine) \\
*n & /n/ & n & *naːpeːwa → \textit{nap} (man) \\
*s & /s/ & s, ss & *soːkya → \textit{sohki} (hard) \\
*š & /ʃ/ & sh & *eːšpowaːki → \textit{ashpowau} (above) \\
*θ & /s/ & s & *eːθkweːwa → \textit{squa} (woman) \\
*l & /n/ & n & *laːkeːwa → \textit{noh} (father) \\
*w & /w/ & w & *waːposwa → \textit{womp} (white) \\
*y & /j/ & y & *ayaːpeːwa → \textit{ayim} (man?) \\
*h & /h/ & h, (unwritten) & *aːhšoːwa → \textit{ashsh} (?) \\
\bottomrule
\end{tabular}
\caption{PA to Massachusett regular correspondences}
\label{tab:1363-pa}
\end{table}

\subsubsection{Confirmed Cognate Mappings for Pipeline Testing}

\begin{table}[h]
\centering
\begin{tabular}{llll}
\toprule
Gloss & Trumbull/Eliot (expected) & Confidence & PA Source \\
\midrule
`no, not' & \textit{matta} & High & PA *mati- \\
`woman' & \textit{squa} & High & PA *eːθkweːwa \\
`man' & \textit{nap} & High & PA *naːpeːwa \\
`house' & \textit{wetu} & High & PA *wiːkiwaːmi \\
`water' & \textit{nippe} & High & PA *nepyi \\
`stone' & \textit{qussuk} & High & PA *aːsenya \\
`fire' & \textit{nuta} (?) & Moderate & PA *eškw- \\
`three' & \textit{ash} & High & PA *nšwaːśika \\
`white' & \textit{womp(ish)} & High & PA *waːposwa (rabbit → white) \\
`devil' & \textit{achm∞wonk} & Moderate & See §\ref{tab:1363-devil} \\
`snake' & \textit{askkuhnk} (?) & Moderate & PA *ašk-? \\
\bottomrule
\end{tabular}
\caption{Confirmed cognate mappings}
\label{tab:1363-cognates}
\end{table}

\subsection{Conversion Workflow}
\texttt{[Pending implementation --- see Issue \#1363]}

\subsection{Linguist Feedback}
\texttt{[Template --- content pending review]}

\subsection{Related Work: Colonial Recorder Perception Problem}

The core challenge of this research card --- English-speaking recorders
misperceiving and mis-transcribing Algonquian phonemes --- is a known
phenomenon documented across Eastern North America. Two researchers' work is
directly relevant:

\textbf{Dr. Keith Cunningham} (Linguist, Nanticoke Indian Tribe). His doctoral
dissertation \emph{A Phonological Analysis of Nanticoke With Practical
Applications for Language Revitalization} (Georgetown University) analyzes
three 18th-century Nanticoke word lists (Heckewelder 1785) using the same
methodology: historical phonology from colonial recorders, with the same
challenges of multilingual informants and orthographic complexity. His 2024
Algonquian Conference presentation \emph{A Revised Phonological Analysis of
the Heckewelder Vocabulary of Nanticoke} directly parallels the Trumbull/Eliot
analysis --- both involve working with standardized missionary orthographies
(Eliot for Massachusett, Heckewelder for Nanticoke) and both require PA
comparative calibration. Cunningham's work demonstrates that the colonial
recorder perception problem extends beyond SNEA into the broader Eastern
Algonquian family.

\textbf{Dr. Craig Kopris} (Independent scholar). His paper \emph{Les sons
wyandots perçus par des oreilles étrangères} (Wyandot sounds perceived by
foreign ears, \emph{Recherches amérindiennes au Québec}, 2014) is the exact
framing of the Trumbull/Eliot problem: how European recorders systematically
misperceived and mis-transcribed indigenous phonemes. His \emph{Wyandot
Phonology: Recovering the Sound System of an Extinct Language} applies the
same recovery methodology to an unrelated language family (Iroquoian),
demonstrating that the foreign-ear perception problem is universal across
colonial documentation of North American languages. His Dictionary of Wyandot
project (NEH award FN-255576-17) shows the long-term application of these
methods to language revitalization.

The Trumbull corpus (4,491 records, the largest in the SNEA collection)
benefits from Eliot's relatively systematic orthography --- but the
preaspiration gap and vowel ambiguity problems documented in this section are
the same foreign-ear perception issues that Cunningham and Kopris address in
their respective language families.

\subsection{Open Questions for Linguist Review}

\begin{enumerate}
  \item \textbf{∞(ꝏ) phonetics:} Is the \texttt{oozzz} mapping phonetically
    accurate, or should it be revised to \texttt{uu} with a diacritic for
    preaspiration? The \texttt{zzz} component seems ad-hoc.
  \item \textbf{Preaspiration inference rules:} Can a reliable set of
    morphotactic rules predict when preaspiration /h/ was present but unwritten
    in Eliot's orthography? Or does this require per-entry comparative
    verification?
  \item \textbf{Dialect stratification:} How much of the variation between
    Trumbull/Eliot (Natick) and Winslow/Wood (Pokanoket) is orthographic vs.
    phonological vs. dialectal? The \textit{achm∞wonk} vs. \textit{hobbamock}
    divergence suggests real dialect differentiation.
  \item \textbf{Apostrophe resolution:} Is there a predictable pattern for when
    \texttt{'} = /h/ vs. /ʔ/ vs. simple vowel elision? Edwards (Mahican)
    apostrophe usage may provide clues.
  \item \textbf{Eliot's \textit{c} before \textit{e}:} Are there identifiable
    exceptions to the \textit{c} before front vowels → /s/ rule?
  \item \textbf{Edwards cross-over:} How much does Edwards' (1787) Mahican
    preaspiration system overlap with Massachusett? If the systems are
    parallel, Edwards' data can directly inform Trumbull /h/ inference.
\end{enumerate}

\subsection{Bibliography}

\begin{enumerate}

\item Aubin, George F. (1975). \emph{A Proto-Algonquian Dictionary.} National
Museum of Man, Mercury Series, Canadian Ethnology Service Paper No. 39.
Ottawa: National Museums of Canada.

\item Costa, David J. (2007). ``The Dialectology of Southern New England
Algonquian.'' In H.C. Wolfart, ed., \emph{Papers of the 38th Algonquian
Conference}, pp. 81--127. Winnipeg: University of Manitoba.

\item Eliot, John. (1663). \emph{Mamusse Wunneetupanatamwe Up-Biblum God.}
Cambridge, Massachusetts: Samuel Green and Marmaduke Johnson.

\item Eliot, John. (1666). \emph{The Indian Grammar Begun.} Cambridge,
Massachusetts.

\item Goddard, Ives, and Kathleen J. Bragdon. (1988). \emph{Native Writings
in Massachusett.} 2 vols. Memoirs of the American Philosophical Society 185.
Philadelphia: American Philosophical Society.

\item O'Brien, Frank Waobu. \emph{Appendix: Guide to Historical Spellings
and Sounds in New England Algonquian Languages.} Aquidneck Indian Council.

\item Trumbull, James Hammond. (1903). \emph{Natick Dictionary.} Bureau of
American Ethnology Bulletin 25. Washington: Smithsonian Institution.

\end{enumerate}
